\chapter{Moldavia: extracción de datos gubernamentales}
\label{chap:incidente3}

\section{Contexto y justificación de selección}
Entre el 19 y el 23 de octubre de 2026, un ataque contra servicios públicos digitales de Moldavia permitió extraer datos personales. El incidente se comunicó públicamente el día 24~\cite{age2026inicial,age2026datos}. Se selecciona por sus consecuencias para la privacidad y la prestación de servicios públicos, y porque permite examinar los riesgos de conectar numerosos registros mediante aplicaciones con controles de acceso deficientes.

\section{Desarrollo del incidente y CVE/CVSS}
Según AGE, los atacantes explotaron servicios accesibles
sin autenticación y utilizaron la información obtenida
para realizar extracciones adicionales~\cite{age2026datos}.
La Presidencia describió fallos de autenticación y control
de acceso que permitieron consultar registros de al menos
diez instituciones. El ataque comenzó la tarde del 19;
el día siguiente se registró aproximadamente un millón
de solicitudes automatizadas, treinta veces el volumen
habitual~\cite{presidencia2026moldavia}.
Las fuentes no publican el código vulnerable ni permiten
identificar con certeza una técnica concreta, como
inyección SQL.

Los comunicados consultados no proporcionan CVE ni CVSS.
Como ejercicio propio, se estima un escenario de lectura
remota no autorizada siguiendo CVSS v3.1~\cite{firstcvss31}:\par
\noindent{\small\ttfamily
    CVSS:3.1/AV:N/AC:L/PR:N/UI:N/S:U/C:H/I:N/A:N}\par
\noindent El resultado es 7,5, alto, para este escenario hipotético de lectura no autorizada. Se asumen complejidad baja y alcance sin cambios; integridad y disponibilidad se fijan en cero para delimitar el cálculo, no porque se haya demostrado la ausencia de esos impactos. La documentación pública resulta insuficiente para asignar una puntuación definitiva a la vulnerabilidad real.

\section{Impacto y mitigación}
AGE estimó inicialmente extracciones asociadas a unos
750.000 identificadores personales. Entre las categorías
expuestas figuraban identificación, relaciones familiares,
prestaciones sociales e información financiera; no todas
afectaban a cada persona~\cite{age2026datos}.
Una evaluación posterior redujo el número de ciudadanos
afectados a aproximadamente 450.000, que debían ser
notificados~\cite{presidencia2026moldavia}.
La información facilita suplantaciones y fraudes dirigidos,
aunque conocer un IBAN no permite por sí solo disponer
del dinero. AGE no identificó modificaciones de los
registros ni exposición de datos de tarjetas
bancarias~\cite{age2026datos}.

La desconexión preventiva dejó servicios temporalmente
inaccesibles~\cite{age2026inicial}.
La respuesta posterior incluyó correcciones en MPass,
verificaciones adicionales de autorización y refuerzo
de las plataformas~\cite{presidencia2026moldavia}.
Como prevención se proponen controles de acceso por
operación y registro, límites de consultas y alertas
ante extracciones anómalas. Para los ciudadanos, AGE
recomendó verificar las solicitudes por canales oficiales
y no facilitar contraseñas ni códigos~\cite{age2026datos}.

\section{Análisis crítico}
El caso evidencia que contener el acceso no revierte
la pérdida de confidencialidad. La recuperación debe
combinar restauración segura, notificación y prevención
del fraude posterior.
También exige distinguir solicitudes, registros e individuos:
sus cifras no son equivalentes. La revisión del número
de afectados demuestra por qué deben fecharse las
estimaciones y actualizarse las conclusiones conforme
avanza la investigación.